% !PRINTED 130
\weberpage{130}
%<<CONT
\gesperrt{theiler von $\mu$, so giebt es $\nu$ und nicht mehr Elemente in $P$, deren Grad in $\nu$ aufgeht.}

Wenden wir diesen Satz auf $\nu = t$ an, so folgt, dass es ausser $T$ keine Elemente in $P$ geben kann, deren Grad $= t$ ist. Wenn aber $c$ irgend ein Element von $P$ ist, so ist die mit $T$ conjugirte Gruppe $c^{-1}Tc$ gleichfalls vom Grade $t$, und muss also mit $T$ identisch sein. Damit ist dann bewiesen, dass $T$ ein Normaltheiler von $P$ ist.

Zum Beweis des Satzes $\alpha$) wendet man die vollständige Induction an. Derselbe ist offenbar richtig, wenn $\mu = 1$, $\nu = n$ ist, und wir setzen also als bewiesen voraus:

$\alpha'$) \gesperrt{Ist $\mu'\nu'$ durch kein Quadrat theilbar, jeder Primtheiler von $\nu'$ grösser als jeder Primtheiler von $\mu'$, und die Anzahl der Primtheiler von $\mu'$ kleiner als die Anzahl der Primtheiler von $\mu$, so giebt es in jeder Gruppe vom Grade $\mu'\nu'$ genau $\nu'$ Elemente, deren Grad ein Theiler von $\nu'$ ist.}

Um daraus den Beweis des Satzes $\alpha$) abzuleiten, bezeichnen wir mit $p$ den grössten Primtheiler von $\mu$ und mit $Q$ eine in $P$ enthaltene Gruppe vom Grade $p$, deren Existenz nach dem \gesperrt{Cauchy-Sylow'schen} Theorem feststeht. Dann sind alle Primfactoren von $\nu$ grösser als $p$. Wie in \S.\,30 bezeichnen wir mit $R$ die Gruppe, die aus allen der Bedingung
\[
c^{-1}Qc = Q
\]
genügenden Elementen $c$ von $P$ besteht, deren Grad wir, wie oben, mit $pr$ bezeichnen; $j$ bedeutet, wie früher, den Index des Theilers $R$ von $P$. Wir setzen $r = r_1 r_2$, und verstehen unter $r_1$ den grössten gemeinschaftlichen Theiler von $\mu$ und $r$, so dass $r_2$ der grösste gemeinschaftliche Theiler von $\nu$ und $r$ ist. Es ist dann
\[
n = \mu\nu = r_1 p r_2 j.
\]
Nach der Hypothese $\alpha'$) enthält nun $P$ genau $p\nu$ Elemente, deren Grad ein Theiler von $p\nu$ ist; es möge mit $U$ die Gesammtheit dieser Elemente bezeichnet sein. Die Gruppe $R$ enthält aber, gleichfalls nach $\alpha'$), genau $pr_2$ Elemente, deren Grad ein Theiler von $pr_2$, also ein Theiler von $p\nu$ ist. Dieses System wollen wir mit $V$ bezeichnen. Offenbar sind alle Elemente von $V$ zugleich in $U$ enthalten.
